\begin{abstract}
Model merging in parameter space has emerged as a powerful paradigm for multi-task learning without retraining. While dynamic test-time ensembling (routing) offers adaptive multi-task capabilities, it incurs prohibitive computational and memory-bandwidth overhead during inference by dynamically reconstructing full model weights for each input batch. This makes dynamic ensembling highly impractical for real-world, latency-sensitive deployments and resource-constrained edge devices. In this paper, we present \textbf{Hybrid-Router}, a highly practical and low-latency hybrid dynamic model merging framework that resolves this bottleneck. Recognizing that early layers in deep networks (such as Vision Transformers) function as task-agnostic feature extractors while late layers capture task-specific specialized representations, we partition the model layer-wise. We statically merge the early layers offline (zero runtime overhead) and dynamically route and merge only the final $k$ layers at test-time. We utilize standard competitive Softmax routers as the primary high-accuracy architecture, while conducting an exploratory study on a Softmax-free independent sigmoidal projection engine (\textbf{BSigmoid-Router}) to analyze the trade-offs and scaling dynamics of uncoupled task activations.

Our evaluations within a controlled Parameter-Space Representation Sandbox modeling high-conflict vision domains (MNIST, FashionMNIST, CIFAR-10, and SVHN) demonstrate that Hybrid-Router forms a highly favorable systems-level Pareto frontier. At $k=4$, our Hybrid-Router achieves a joint mean of \textbf{76.75\%} (a massive \textbf{+4.44\%} improvement over the offline static AdaMerging baseline of 72.31\%) while cutting weight reconstruction latency and task-vector VRAM footprint by \textbf{71.3\%} (2.95 ms vs 10.28 ms) and saving 71.4\% of active dynamic task-vector storage. Furthermore, under highly constrained calibration splits, we find that restricting the learnable space of the routing optimizer via layer partitioning acts as a form of structural regularization, yielding an accuracy of \textbf{84.79\%} at $k=12$ (a minor \textbf{+0.22\%} improvement over fully dynamic ensembling at $k=14$) while saving \textbf{14.3\%} in latency. Under streaming noise at batch size $B=1$, our regularized routing model maintains a robust stream accuracy of \textbf{84.55\%}. This work demonstrates that we can achieve exceptional multi-task capabilities with a highly favorable, customizable latency-accuracy-memory trade-off.
\end{abstract}
